\noindent\textbf{Role tréninkových dat (stručně)} \emph{(general background knowledge)}: Tréninková data určují vzory, fakta a možná zkreslení v odpovědích; kvalita, rozmanitost a časové pokrytí ovlivňují přesnost a aktuálnost.\\[6pt]
\noindent\textbf{Co znamená cut‑off dat} \emph{(general background knowledge)}: Cut‑off je poslední časový bod zahrnutí tréninkových dat — model nemusí znát události po tomto datu; v učebním kontextu zdůrazněte ověřování z externích zdrojů.\\[6pt]
\noindent\textbf{Klíčové pedagogické body k zahrnutí do schématu} \begin{itemize}
  \item Zdroj dat: weby, knihy, interní dokumenty, anotace; kvalita a bias: reprezentativnost, chyby, systematická zkreslení.
  \item Časová osa / cut‑off: vizuální osa znázorňující „do kdy“ model zná informace.
  \item Soukromí a právní aspekty: osobní údaje, DPIA; Retrieval/RAG: jak indexy a retrieval vrstva mění chování.
  \item Didaktika: kdy vyžadovat ověření a ukázat zdroje.
\end{itemize}
\noindent\textbf{Návrh jednoduchého schématu pro slide (1–2 snímky)} \emph{(general background knowledge)}: \begin{enumerate}
  \item Snímek 1 — „Mapa tréninkových dat“: centrální kruh „Model“ obklopený bublinami „Web“, „Knihy“, „Interní dokumenty“; pod tím osa s cut‑off.
  \item Snímek 2 — „Dopady a mitigace“: tři sloupce: rizika, příznaky v odpovědích, opatření (ověřit zdroj, DPIA, lidská validace).
\end{enumerate}
\noindent\textbf{Jak toto schéma prezentovat bez matematiky} \emph{(general background knowledge)}: \begin{itemize}
  \item Metafora „knihovny“ a krátké demonstrace: porovnejte otázky o událostech před a po cut‑off.
  \item Praktická pravidla: vždy ukázat zdroj a provést lidskou kontrolu u kritických tvrzení.
\end{itemize}
\noindent\textbf{Krátký workshop / aktivita ve třídě (praktický příklad)}: \begin{enumerate}
  \item Fáze A (bez AI): studenti vypracují krátký text.
  \item Fáze B (s AI): studenti použijí model k rozšíření textu a porovnají přínosy a chyby; diskutujte odkud model čerpá informace \cite{kb_malinka_chatgpt_ve_vyuce_dva_roky_pote_2024_pdf_90bc3aaf_page_12_1}.
  \item Reflexe a alternativa: studenti popíší ověřování tvrzení; jako alternativa použijte didaktický svět (např. Minecraft Education) k názornému zobrazení vlivu tréninkových dat \cite{kb_ai_101_tipu_pdf_04e4a115_page_15_2}.
\end{enumerate}
\noindent\textbf{Krátké bezpečnostní a etické připomenutí} \cite{kb_malinka_chatgpt_ve_vyuce_dva_roky_pote_2024_pdf_90bc3aaf_page_12_1}: zvažte dopad na ochranu soukromí a zařaďte analýzu rizik (DPIA) při práci s osobními nebo interními daty.\\[6pt]
\noindent\textbf{Doporučené diskusní otázky pro studenty} \emph{(general background knowledge)}: \begin{itemize}
  \item Jaké zdroje byste považovali za důvěryhodné pro ověření výroku z modelu?
  \item Co by model mohl „nevědět“ kvůli cut‑off datu a jak to rozpoznat?
  \item Jak mitigovat riziko, že model předá citlivé informace z tréninkových dat?
\end{itemize}